Se caracteriza un filtro pasabajos RC de primer orden ($R = 1{,}6\,\mathrm{k}\Omega$, $C = 100\,\mathrm{nF}$), contrastando tres fuentes de datos: la respuesta \textbf{teórica} obtenida del modelo analítico, la \textbf{simulada} con ngspice y la \textbf{medida} sobre el circuito armado.

La frecuencia de corte nominal es
\[ f_c = \frac{1}{2\pi RC} = \frac{1}{2\pi \cdot 1600 \cdot 100\times10^{-9}} \approx 994{,}7\ \mathrm{Hz}. \]
